\chapter{A Complete Example: A CORDIC Rotator}
\label{ch:case}

\chapabstract{One design of real, if modest, complexity, written twice. This
chapter contains the complete \VHDL{} and \SV{} sources for an iterative
CORDIC rotator, four testbenches for it in three languages, a Python golden
model, and the evidence that all of them agree. Every file is reproduced in
full, and every line of output quoted here was produced by actually running
the tools.}

\section{What We Are Building, and Why This Example}
\label{sec:case:what}

\index{CORDIC}The block is \code{cordic\_rot}: a rotation-mode CORDIC that
takes a vector $(x, y)$ and an angle $\theta$ and returns the vector rotated
by $\theta$, scaled by a fixed gain. Feed it $x = \text{\code{X0\_UNIT}}$,
$y = 0$ and it hands you $(\cos\theta, \sin\theta)$.

That is a small specification. What makes it the right example for this book
is everything the implementation needs in order to satisfy it:

\begin{itemize}
  \item a \textbf{fixed-point number format} that somebody has to choose, and
        that neither language will check for you in the same way;
  \item a \textbf{constant table} computed at elaboration from a
        transcendental function, which is where \VHDL{}'s
        \code{ieee.math\_real} and \SV{}'s \code{\$atan} meet;
  \item a \textbf{control state machine}, small enough to read on one page;
  \item a \textbf{parameterised datapath} with shifts, adds and a saturating
        output stage;
  \item a \textbf{handshake protocol} on both ports, which is what makes a
        proper monitor and scoreboard necessary rather than decorative;
  \item and \textbf{arithmetic} whose correct answer is not obvious by
        inspection, so a golden model earns its keep.
\end{itemize}

If you understand every line of this chapter, you can write \SV{}. There is
nothing else in the language you will meet in your first year that is not a
variation on something here.

\Cref{fig:case:block} is the whole block. Three stages and a controller: a
coarse rotation that folds the angle into the range where the algorithm
converges, a micro-rotation datapath that iterates $N$ times, and a saturating
narrow-down on the way out.

\begin{figure}[htbp]
  \centering
  \begin{tikzpicture}[
      font=\small\sffamily,
      node distance=6mm,
      blk/.style={draw=codekey,fill=svbg,line width=0.9pt,rounded corners=2pt,
                  minimum height=11mm,minimum width=26mm,align=center,
                  inner sep=4pt},
      rom/.style={draw=accent,fill=tipbg,line width=0.9pt,rounded corners=2pt,
                  minimum height=9mm,minimum width=20mm,align=center},
      ctl/.style={draw=accentalt,fill=warnbg,line width=0.9pt,
                  rounded corners=2pt,minimum height=9mm,minimum width=26mm,
                  align=center},
      ar/.style={-{Latex[length=2.2mm]},draw=codenumber,thick}]

    \node[blk] (coarse) {coarse\\rotation};
    \node[blk,right=14mm of coarse] (micro) {micro-rotation\\$x,y,z$};
    \node[blk,right=14mm of micro] (sat) {saturating\\narrow};
    \node[rom,below=9mm of micro] (romn) {\code{ATAN\_ROM}};
    \node[ctl,above=9mm of micro] (fsm) {FSM\\ \scriptsize IDLE / ROTATE / DONE};

    \node[left=11mm of coarse,align=right,font=\scriptsize\ttfamily]
      (in) {req\_x\\req\_y\\req\_theta};
    \node[right=11mm of sat,align=left,font=\scriptsize\ttfamily]
      (out) {rsp\_x\\rsp\_y};

    \draw[ar] (in.east)     -- (coarse.west);
    \draw[ar] (coarse.east) -- (micro.west);
    \draw[ar] (micro.east)  -- (sat.west);
    \draw[ar] (sat.east)    -- (out.west);
    \draw[ar] (romn.north)  -- (micro.south);
    \draw[ar] (fsm.south)   -- (micro.north);

    % the fold-back path: one iteration per clock
    \draw[ar] ([yshift=3mm]micro.east) -- ++(0.7,0) -- ++(0,-1.35)
              -- ++(-4.6,0) -- ([yshift=-3mm]micro.west);
    \node[font=\scriptsize,below=1.5mm of micro,xshift=15mm,yshift=-1mm]
      {one iteration per clock};

    \node[font=\scriptsize\itshape,left=1mm of coarse,yshift=-9mm,
          anchor=north east] {\phantom{x}};
    \node[font=\scriptsize,above=0.5mm of in] {\itshape valid / ready};
    \node[font=\scriptsize,above=0.5mm of out] {\itshape valid / ready};
  \end{tikzpicture}
  \caption{The \code{cordic\_rot} datapath. The coarse-rotation stage runs
  once per transaction; the micro-rotation stage runs $N$ times, reusing one
  adder pair and one shifter pair. A transaction takes $N + 2$ clock cycles.}
  \label{fig:case:block}
\end{figure}

\begin{notebox}{Folded, not pipelined}
This is the area-efficient shape: one set of adders, reused $N$ times, at a
throughput of one transaction every $N + 2$ cycles. The other shape unrolls
the loop into $N$ pipeline stages and accepts one transaction per cycle at
roughly $N$ times the area. \Cref{sec:case:exercises} asks you to build that
version, because the interesting part is not the datapath but what happens to
the handshake.
\end{notebox}

\section{The Algorithm}
\label{sec:case:algorithm}

\index{CORDIC!rotation mode}A rotation by $\theta$ is
\[
  \begin{pmatrix} x' \\ y' \end{pmatrix} =
  \begin{pmatrix} \cos\theta & -\sin\theta \\
                  \sin\theta & \cos\theta \end{pmatrix}
  \begin{pmatrix} x \\ y \end{pmatrix} ,
\]
which needs four multiplications and two transcendental functions. Volder's
observation in 1959 was that if you restrict yourself to rotations by the
particular angles $\alpha_i = \arctan 2^{-i}$, the matrix becomes
\[
  \begin{pmatrix} x_{i+1} \\ y_{i+1} \end{pmatrix} =
  \cos\alpha_i
  \begin{pmatrix} 1 & -d_i 2^{-i} \\ d_i 2^{-i} & 1 \end{pmatrix}
  \begin{pmatrix} x_i \\ y_i \end{pmatrix} ,
  \qquad d_i = \pm 1 ,
\]
and the only operations left inside the matrix are a shift and an add. Drop
the $\cos\alpha_i$ scale factor, keep a running account of how much angle is
left to rotate, and you have the algorithm:

\begin{align}
  x_{i+1} &= x_i - d_i \, (y_i \gg i) \label{eq:case:x}\\
  y_{i+1} &= y_i + d_i \, (x_i \gg i) \label{eq:case:y}\\
  z_{i+1} &= z_i - d_i \, \alpha_i    \label{eq:case:z}
\end{align}

with $z_0 = \theta$ and $d_i = +1$ if $z_i \ge 0$, $-1$ otherwise. The
decision is the sign bit of $z_i$: no comparison, no subtraction, one wire.

Three consequences follow, and all three appear in the code.

\subsection{The processing gain}

Dropping $\cos\alpha_i$ at every step means the vector grows. After $N$
iterations the length has been multiplied by
\[
  K_N = \prod_{i=0}^{N-1} \sqrt{1 + 2^{-2i}} ,
\]
which converges very quickly: $K_{14} = 1.6467602540$ and the limit is
$K_\infty = 1.6467602581$. The gain is a constant, so you can cancel it once,
in advance, by pre-scaling the input. The package exposes
\code{X0\_UNIT} $= \operatorname{round}(2^{14}/K) = 9949$ for exactly this
purpose.

\begin{notebox}{Which $K$?}
The testbenches use $K = 1.6467602581210656$, the limit, while the true gain
for $N = 14$ is $1.6467602540312922$. The difference is $4 \times 10^{-9}$,
about $6.7\times10^{-5}$ of one LSB in Q1.14, so it is invisible next to the
$2.5$ LSB of real quantisation error. Using the limit is the conventional
choice because it does not change when you change $N$. If you ever tighten
the tolerance below one LSB, use the exact product for your $N$.
\end{notebox}

\subsection{The convergence range}

The angle left over after all $N$ steps can only be driven to zero if the
step you are about to take is never larger than the sum of all the steps
remaining. That gives
\[
  |\theta| \le \sum_{i=0}^{\infty} \arctan 2^{-i} = 1.7433 \text{ rad}
  \approx 99.9^\circ .
\]
So the raw recurrence handles a little under $\pm100^\circ$, and nothing more.

\subsection{Quadrant folding}

To cover the full circle you rotate by a whole quadrant first, which costs
nothing because a rotation by $\pm\pi/2$ is a swap and a negation:
\[
  +\tfrac{\pi}{2}: (x,y) \mapsto (-y, x) , \qquad
  -\tfrac{\pi}{2}: (x,y) \mapsto (y, -x) .
\]
One fold brings $|\theta| \le \pi$ down to $|\theta| \le \pi/2 = 1.5708$,
comfortably inside $1.7433$. Two folds would cover $|\theta| \le 3\pi/2$, and
the design does not do them: it accepts $|\theta| \le \pi$ and fires an
assertion otherwise. Range reduction of an arbitrary angle belongs upstream,
where it can be done once for a whole stream.

\Cref{fig:case:converge} shows what the micro-rotations actually do. Each one
overshoots, and the overshoot halves.

\begin{figure}[htbp]
  \centering
  \begin{tikzpicture}[font=\small\sffamily,scale=1.0]
    \def\R{3.0}
    % unit circle
    \draw[codeframe,line width=0.8pt] (0,0) circle (\R);
    \draw[codenumber,-{Latex[length=1.8mm]}] (-\R-0.4,0) -- (\R+0.5,0)
      node[right,font=\scriptsize]{$x$};
    \draw[codenumber,-{Latex[length=1.8mm]}] (0,-0.4) -- (0,\R+0.5)
      node[above,font=\scriptsize]{$y$};

    % target angle 60 degrees
    \draw[accentalt,line width=1.3pt,-{Latex[length=2.4mm]}]
      (0,0) -- (60:\R) node[above right,font=\scriptsize]{target $\theta$};
    \draw[accentalt,dashed] (0,0) -- (60:\R+0.45);

    % successive approximations: 45, 71.6, 45+26.6-14.0=58.3, +7.1=65.4
    \draw[codekey,line width=1.0pt,-{Latex[length=2mm]}] (0,0) -- (0:\R);
    \node[font=\scriptsize,codekey] at (2:\R+0.32) {$i{=}0$};
    \draw[codekey,line width=1.0pt,-{Latex[length=2mm]}] (0,0) -- (45:\R);
    \node[font=\scriptsize,codekey] at (45:\R+0.32) {$1$};
    \draw[codekey,line width=1.0pt,-{Latex[length=2mm]}] (0,0) -- (71.6:\R);
    \node[font=\scriptsize,codekey] at (73:\R+0.3) {$2$};
    \draw[codekey,line width=1.0pt,-{Latex[length=2mm]}] (0,0) -- (57.5:\R);
    \node[font=\scriptsize,codekey] at (55:\R+0.32) {$3$};
    \draw[codekey,line width=1.0pt,-{Latex[length=2mm]}] (0,0) -- (64.6:\R);
    \node[font=\scriptsize,codekey] at (65.5:\R+0.3) {$4$};

    \node[align=left,font=\scriptsize,anchor=north west]
      at (-\R-0.3,-0.6)
      {Each step rotates by $\arctan 2^{-i}$: $45.0^\circ$, $26.6^\circ$,\\
       $14.0^\circ$, $7.1^\circ$, \ldots\ The sign alternates as the\\
       residual angle $z$ changes sign. After 14 steps the\\
       residual is below the resolution of the angle format.};
  \end{tikzpicture}
  \caption{The first four micro-rotations converging on a target of
  $60^\circ$. The vector also grows by the processing gain at every step;
  that growth is not drawn.}
  \label{fig:case:converge}
\end{figure}

The table of angles is fixed at elaboration. \Cref{tab:case:atan} is the one
this design uses, taken verbatim from the generated file
\file{sv/cordic\_atan\_rom.svh}.

\begin{table}[htbp]
  \centering
  \small
  \caption{The arctangent ROM. Values are in the internal angle format,
  Q2.15, so one unit is $2^{-15} = 3.05\times10^{-5}$ rad.}
  \label{tab:case:atan}
  \begin{tabular}{@{}rrr@{\hspace{2.5em}}rrr@{}}
    \toprule
    $i$ & $\arctan 2^{-i}$ (rad) & Q2.15 &
    $i$ & $\arctan 2^{-i}$ (rad) & Q2.15\\
    \midrule
     0 & 0.785398163 & 25736 &  7 & 0.007812341 &  256\\
     1 & 0.463647609 & 15193 &  8 & 0.003906230 &  128\\
     2 & 0.244978663 &  8027 &  9 & 0.001953123 &   64\\
     3 & 0.124354995 &  4075 & 10 & 0.000976562 &   32\\
     4 & 0.062418810 &  2045 & 11 & 0.000488281 &   16\\
     5 & 0.031239833 &  1024 & 12 & 0.000244141 &    8\\
     6 & 0.015623729 &   512 & 13 & 0.000122070 &    4\\
    \bottomrule
  \end{tabular}
\end{table}

\begin{notebox}{Where the table stops paying}
From $i = 5$ onwards the entries are exactly $2^{15-i}$: the arctangent has
become indistinguishable from its argument at this precision. That is the
signal that $N = 14$ is matched to a 15-bit fractional angle, and that adding
a fifteenth iteration would buy nothing without widening the angle first.
\end{notebox}

\section{Choosing the Number Format}
\label{sec:case:format}

This is the section where a \VHDL{} designer discovers what \SV{} will not do
for them.

\subsection{The decisions}

\textbf{$x$ and $y$: signed 16 bits, 14 fractional (Q1.14).} The natural
output is a unit vector, so the range needs to cover $[-1, 1]$ with a little
room. One integer bit gives $[-2, 2)$ and a resolution of
$2^{-14} = 6.1\times10^{-5}$. Sixteen bits because that is what fits a
DSP block input and what the rest of the world will hand you.

\textbf{$\theta$: signed 16 bits, 13 fractional (Q2.13).} The angle has to
reach $\pi = 3.14159$, so two integer bits are the minimum; the range is
$[-4, 4)$ and the resolution is $1.22\times10^{-4}$ rad. Note that this is
\emph{coarser} than the output resolution: an angle error of
$1.22\times10^{-4}$ rad produces an output error of up to $1.22\times10^{-4}$,
which is two LSB of Q1.14. The input quantisation, not the algorithm, is the
dominant error term.

\textbf{Guard bits: two at each end.} \code{GL = 2} extra fractional bits
below the LSB stop the fourteen successive truncations in
\cref{eq:case:x,eq:case:y} from accumulating into the result; each shift
throws away up to one LSB, and fourteen of them would otherwise cost about
four. \code{GM = 2} extra bits above the MSB give the datapath room for the
processing gain: an input of magnitude $2$ becomes $3.3$ inside the loop, and
without headroom it would wrap.

\textbf{The output saturates rather than wraps.} A CORDIC that silently
returns $-1.99$ when the answer was $+2.01$ is worse than one that returns
$+1.99$, because the wrap looks like a sign error and takes an afternoon to
find.

\Cref{tab:case:formats} collects every constant.

\begin{table}[htbp]
  \centering
  \small
  \caption{Every number format in the design. The internal formats are the
  external ones widened by \code{GL} at the bottom and \code{GM} at the top.}
  \label{tab:case:formats}
  \begin{tabular}{@{}llrrll@{}}
    \toprule
    Quantity & Type & Bits & Frac. & Range & Resolution\\
    \midrule
    $x$, $y$ in/out    & \code{data\_t}   & 16 & 14 & $[-2, 2)$ & $6.1\times10^{-5}$\\
    $\theta$ in        & \code{angle\_t}  & 16 & 13 & $[-4, 4)$ & $1.22\times10^{-4}$\\
    $x$, $y$ internal  & \code{idata\_t}  & 20 & 16 & $[-8, 8)$ & $1.5\times10^{-5}$\\
    $z$ internal       & \code{iangle\_t} & 18 & 15 & $[-4, 4)$ & $3.05\times10^{-5}$\\
    \midrule
    \multicolumn{6}{@{}l}{\itshape Derived constants}\\
    \code{PI\_EXT}      & \multicolumn{2}{l}{$\pi \cdot 2^{13}$}   & \multicolumn{3}{l}{25736 \quad the accepted input limit}\\
    \code{HALF\_PI\_INT}& \multicolumn{2}{l}{$\frac{\pi}{2} \cdot 2^{15}$} & \multicolumn{3}{l}{51472 \quad the coarse-rotation threshold}\\
    \code{X0\_UNIT}     & \multicolumn{2}{l}{$2^{14}/K$}           & \multicolumn{3}{l}{9949 \quad\ the unit-circle seed}\\
    \code{MAX\_DATA}    & \multicolumn{2}{l}{$2^{15}-1$}           & \multicolumn{3}{l}{32767 \quad the saturation limit}\\
    \bottomrule
  \end{tabular}
\end{table}

\subsection{The error budget}

Three terms, and it is worth being able to write them down before you
simulate, so that a measured error you did not predict is a signal rather
than a shrug.

\begin{enumerate}
  \item \textbf{Angle quantisation.} $\pm\frac12$ LSB of Q2.13, so
        $6.1\times10^{-5}$ rad, giving up to $1.0$ LSB of Q1.14 output error.
  \item \textbf{Residual angle after $N$ steps.} Bounded by $\alpha_{N-1} =
        1.22\times10^{-4}$ rad, another $2.0$ LSB, but in practice much less
        because the residual is roughly uniform in
        $[-\alpha_{N-1}, \alpha_{N-1}]$.
  \item \textbf{Truncation in the shifts.} Fourteen right shifts, each losing
        up to one internal LSB of $2^{-16}$, and the guard bits mean those
        internal LSBs are a quarter of an output LSB. Call it $1$ LSB.
\end{enumerate}

Sum: under $5$ LSB. Measured, over 308 random and directed transactions:
$2.55$ LSB in the \VHDL{} testbench, $2.82$ in the \SV{} one. The tolerance
in every testbench is set to $12$ LSB, which leaves room for the tails without
being loose enough to hide a wrong sign, a wrong shift or a missing
micro-rotation. All three of those bugs move the answer by hundreds of LSB.

\begin{gotcha}{Nobody is tracking your binary point}
\VHDL{}-2008 ships \code{ieee.fixed\_pkg}. It gives you \code{sfixed(1
downto -14)}, and when you multiply an \code{sfixed(1 downto -14)} by an
\code{sfixed(3 downto -13)} the compiler works out that the result is an
\code{sfixed(5 downto -27)} and tells you so. Shifts, resizes and the
binary-point bookkeeping are all checked at analysis time.

\SV{} has nothing of the kind. \code{data\_t} is twenty bits and the binary
point exists only in \cref{tab:case:formats}, in the comments, and in your
head. Every one of the decisions above is a comment that the tool will never
verify. This is the single largest thing a \VHDL{} designer gives up when
moving to \SV{} for arithmetic work, and the mitigation is entirely
procedural: name the formats in the package, write the widths in the
comments, and check them against a golden model that was written
independently.
\end{gotcha}

\begin{ruleofthumb}
Write the number format down before you write the RTL, put it in the package
header, and have a model that will disagree with you if you get it wrong.
\end{ruleofthumb}

\section{The Golden Model Comes First}
\label{sec:case:model}

\index{golden model}The Python model in \cref{lst:case:model} was written
before either implementation, and both were checked against it.

Two things about it are deliberate.

It is \textbf{bit-exact}, not floating point. A floating point model tells
you the design computes a rotation; a bit-exact model tells you the design
computes \emph{this} rotation, with this rounding, these guard bits and this
saturation. The second catches a class of bug the first cannot see: a shift
that rounds the wrong way, a guard bit dropped one stage too early, a
saturation limit off by one. Both checks appear in this chapter, and they
fail on different things.

It exploits the fact that \textbf{Python's \code{>>} on a negative integer
floors}, which is exactly what an arithmetic shift right does in hardware.
\code{-7 >> 1} is $-4$ in Python and in \VHDL{}'s \code{shift\_right} and in
\SV{}'s \code{>>>}. It is $-3$ in C, which rounds toward zero. Writing the
model in Python rather than C removed a whole category of mismatch before it
could happen.

\begin{gotcha}{A model that copies the RTL is worth nothing}
The temptation, once the RTL exists, is to write the model by reading it.
The result agrees with the RTL perfectly, including on every bug. The model
here was written from \cref{eq:case:x,eq:case:y,eq:case:z} and the format
table, and the first time it was run against the \VHDL{} it disagreed. The
\VHDL{} was wrong. That is the entire value proposition.
\end{gotcha}

Running the model on its own checks it against the ideal rotation:

\begin{txtcode}
$ make model
N = 14, K = 1.646760, 1/K = 0.607253, X0_UNIT = 9949
HALF_PI = 51472, PI = 102944
atan table: [25736, 15193, 8027, 4075, 2045, 1024, 512, 256, 128, 64, 32,
             16, 8, 4]
worst |error| over [-pi, pi] = 0.000233  (3.82 LSB)
\end{txtcode}

$3.82$ LSB over a 721-point sweep of the full circle, against a budget of
$5$. The model is behaving as the arithmetic says it should, which means it
is fit to be the reference.

\clearpage
\pyfile{code/cordic/python/cordic_model.py}{\file{cordic\_model.py} --- the
bit-exact golden model. Every line maps onto one line of RTL.}{lst:case:model}

\section{The VHDL Implementation}
\label{sec:case:vhdl}

Two files: a package that holds the formats, the types and the table, and an
entity that holds the logic. The complete sources are
\cref{lst:case:vhdlpkg,lst:case:vhdlrot} at the end of this section; the
pieces worth stopping on are here.

\subsection{The package}

The subtypes are the first thing a reader meets, and they carry real
information:

\begin{vhdlcode}
subtype data_t   is signed(DW - 1 downto 0);
subtype angle_t  is signed(AW - 1 downto 0);
subtype idata_t  is signed(IW - 1 downto 0);
subtype iangle_t is signed(IA - 1 downto 0);
\end{vhdlcode}

A subtype constrains an existing type. \code{data\_t} and
\code{signed(15 downto 0)} are interchangeable, and every assignment between
them is checked for width and for direction. Assign an \code{idata\_t} to a
\code{data\_t} and the analyser stops you.

The arctangent table is computed by the analyser:

\begin{vhdlcode}
function build_atan_rom return atan_rom_t is
  variable rom : atan_rom_t;
begin
  for i in 0 to N - 1 loop
    rom(i) := to_signed(
                integer(round(arctan(2.0 ** (-i)) * 2.0 ** IQA)), IA);
  end loop;
  return rom;
end function build_atan_rom;
\end{vhdlcode}

\code{arctan} and \code{round} come from \code{ieee.math\_real}, which is a
simulation and elaboration library: the call is evaluated once, when the
constant is elaborated, and only the fourteen integers reach the netlist.
\Cref{tab:case:atan} is literally the output of this function.

The constant itself is \emph{deferred}: the package declares that it exists
and the package body says what it is.

\begin{vhdlcode}
-- in the package declaration
constant ATAN_ROM : atan_rom_t;
-- in the package body, after build_atan_rom
constant ATAN_ROM : atan_rom_t := build_atan_rom;
\end{vhdlcode}

\begin{tipbox}{Use a deferred constant to initialise from a function}
Writing \code{constant ATAN\_ROM : atan\_rom\_t := build\_atan\_rom;} in the
package declaration works, but it calls the function from the declaration
while the body that defines it has not been elaborated, and analysers warn
about that. The deferred form removes the ambiguity and costs one line. \SV{} has no equivalent: a \code{localparam} must be given
its value where it is declared.
\end{tipbox}

The saturating narrow-down is the last piece of arithmetic in the design:

\begin{vhdlcode}
function sat_narrow (v : idata_t) return data_t is
  constant TW : positive := IW - GL;
  variable t  : signed(TW - 1 downto 0);
begin
  t := v(IW - 1 downto GL);          -- drop the guard bits
  if t > to_signed(MAX_DATA, TW) then
    return to_signed(MAX_DATA, DW);
  elsif t < to_signed(MIN_DATA, TW) then
    return to_signed(MIN_DATA, DW);
  else
    return resize(t, DW);
  end if;
end function sat_narrow;
\end{vhdlcode}

Slicing off the bottom \code{GL} bits of a \code{signed} is exactly an
arithmetic shift right, and it rounds toward $-\infty$, matching the Python
model.

\subsection{The architecture}

The coarse rotation is one \code{process (all)}:

\begin{vhdlcode}
x_ext <= shift_left(resize(req_x_i,     IW), GL);
z_ext <= shift_left(resize(req_theta_i, IA), GL);

p_coarse : process (all) is
begin
  if z_ext > HALF_PI_S then
    x_load <= -y_ext;  y_load <=  x_ext;
    z_load <=  z_ext - HALF_PI_S;
  elsif z_ext < -HALF_PI_S then
    x_load <=  y_ext;  y_load <= -x_ext;
    z_load <=  z_ext + HALF_PI_S;
  else
    x_load <= x_ext;   y_load <= y_ext;   z_load <= z_ext;
  end if;
end process p_coarse;
\end{vhdlcode}

\code{resize} on a \code{signed} sign-extends; \code{shift\_left} adds the
guard bits. Both are explicit, both are checked. If \code{IW} and the width
of \code{x\_ext} ever disagree, the analyser says so.

The micro-rotation is four concurrent assignments, which is as close as
\VHDL{} gets to writing down \cref{eq:case:x,eq:case:y,eq:case:z} directly:

\begin{vhdlcode}
ccw    <= not z_q(IA - 1);            -- z >= 0
x_sh   <= shift_right(x_q, iter_q);
y_sh   <= shift_right(y_q, iter_q);
atan_i <= ATAN_ROM(iter_q);

x_rot <= x_q - y_sh   when ccw = '1' else x_q + y_sh;
y_rot <= y_q + x_sh   when ccw = '1' else y_q - x_sh;
z_rot <= z_q - atan_i when ccw = '1' else z_q + atan_i;
\end{vhdlcode}

\code{shift\_right} on a \code{signed} is arithmetic by definition. There is
no way to write the wrong one.

\begin{notebox}{\code{iter\_q} is a constrained integer}
\code{signal iter\_q : natural range 0 to ITER - 1;} declares a counter that
cannot legally hold 14. If a bug ever drove it out of range, the simulation
would stop with a bounds error naming the line. The \SV{} counter is
\code{logic [3:0]} and would wrap to zero in silence.
\end{notebox}

\clearpage
\vhdlfile{code/cordic/vhdl/cordic_pkg.vhd}{\file{cordic\_pkg.vhd} --- formats,
types, the elaboration-time arctangent ROM and the saturating narrow-down.}{lst:case:vhdlpkg}

\clearpage
\vhdlfile{code/cordic/vhdl/cordic_rot.vhd}{\file{cordic\_rot.vhd} --- the VHDL
implementation.}{lst:case:vhdlrot}

\clearpage
\section{The SystemVerilog Implementation}
\label{sec:case:sv}

The same two files, in the same order, doing the same job. Read the
side-by-side pairs below and then the complete sources in
\cref{lst:case:svpkg,lst:case:svrot}.

\subsection{Package and types}

\begin{rosetta}
\begin{rleft}
package cordic_pkg is
  constant DW : positive := 16;
  subtype data_t is
    signed(DW - 1 downto 0);
  ...
end package cordic_pkg;

-- in the design unit:
library work;
  use work.cordic_pkg.all;
\end{rleft}
\begin{rright}
package cordic_pkg;
  parameter int unsigned DW = 16;
  typedef logic signed [DW-1:0]
          data_t;
  ...
endpackage : cordic_pkg

// in the design unit:
module cordic_rot
  import cordic_pkg::*;
\end{rright}
\end{rosetta}

A \code{typedef} creates a genuinely new name, but it carries no constraint:
\code{data\_t} is sixteen bits interpreted as two's complement, and that is
all the tool knows. Assigning a twenty-bit value to it truncates, silently.
The \code{import} on the module header is the tidy form, and it scopes the
import to that module rather than to the compilation unit.

\subsection{The constant table}

\begin{rosetta}
\begin{rleft}
-- declaration
constant ATAN_ROM : atan_rom_t;

-- body
constant ATAN_ROM : atan_rom_t
  := build_atan_rom;
\end{rleft}
\begin{rright}
function automatic atan_rom_t
    build_atan_rom();
  ...
endfunction

localparam atan_rom_t ATAN_ROM
  = build_atan_rom();
\end{rright}
\end{rosetta}

Inside the function, \code{\$atan} plays the part of \code{arctan} and
\code{\$rtoi} the part of \code{integer(round(...))}:

\begin{svcode}
rom[i] = iangle_t'($rtoi($atan(2.0 ** (-i)) * (2.0 ** IQA) + 0.5));
\end{svcode}

\code{\$rtoi} truncates toward zero, so the \code{+ 0.5} is doing the
rounding that \VHDL{}'s \code{round} does for you. Forget it and every entry
in the table is one too small, the rotations are all slightly short, and the
error grows with the angle: a bug that a tolerance of 12 LSB would not catch
at small angles and would catch at large ones.

\begin{gotcha}{A constant function is not portable}
A function that returns an unpacked array, used to initialise a
\code{localparam}, is legal IEEE 1800 and works in \tool{Verilator} 5.020 and
in the commercial simulators. \tool{Icarus Verilog} 12 rejects it outright:
\code{Array rom needs an array index here} followed by \code{Unable to
evaluate parameter ATAN\_ROM}. Some synthesis front ends refuse
\code{\$atan} in a constant expression too. When portability matters, generate
the table from a script instead --- \cref{sec:case:python} does exactly that,
and \file{sv/cordic\_atan\_rom.svh} is the result.
\end{gotcha}

\subsection{Widening and coarse rotation}

\begin{rosetta}
\begin{rleft}
x_ext <= shift_left(
           resize(req_x_i, IW), GL);
z_ext <= shift_left(
           resize(req_theta_i, IA), GL);
\end{rleft}
\begin{rright}
x_ext = idata_t'(req_x_i)      <<< GL;
z_ext = iangle_t'(req_theta_i) <<< GL;
\end{rright}
\end{rosetta}

The cast sign-extends because both types are signed. This is one of the few
places where \SV{}'s implicit width rules genuinely help: the cast states the
intent and the shift does the rest, in one line instead of two nested calls.
It is also one line in which a typo -- \code{iangle\_t} where you meant
\code{idata\_t} -- compiles happily and truncates two bits.

\subsection{The micro-rotation}

\begin{rosetta}
\begin{rleft}
x_sh <= shift_right(x_q, iter_q);

x_rot <= x_q - y_sh
         when ccw = '1'
         else x_q + y_sh;
\end{rleft}
\begin{rright}
x_sh = x_q >>> iter_q;

if (ccw) x_rot = x_q - y_sh;
else     x_rot = x_q + y_sh;
\end{rright}
\end{rosetta}

\begin{gotcha}{\code{>>} and \code{>>>} look almost the same}
\code{x\_q >> iter\_q} compiles, simulates, and gives the right answer for
every positive input. It is wrong for every negative one, and half of a
random stimulus stream is negative, so the failure is loud --- if you have a
scoreboard. If your testbench prints waveforms for a human to inspect, and
the human happens to look at a positive case, this bug ships. \VHDL{} makes
it unwriteable.
\end{gotcha}

\subsection{The state machine}

\begin{rosetta}
\begin{rleft}
type state_t is
  (ST_IDLE, ST_ROTATE, ST_DONE);

p_fsm : process (all) is
begin
  state_d <= state_q;
  req_ready_o <= '0';
  case state_q is
    when ST_IDLE =>
      req_ready_o <= '1';
      ...
  end case;
end process p_fsm;
\end{rleft}
\begin{rright}
typedef enum logic [1:0] {
  ST_IDLE, ST_ROTATE, ST_DONE
} state_e;

always_comb begin
  state_d = state_q;
  req_ready_o = 1'b0;
  unique case (state_q)
    ST_IDLE: begin
      req_ready_o = 1'b1;
      ...
  endcase
end
\end{rright}
\end{rosetta}

Three differences, all of them consequential.

The \VHDL{} enumeration has no encoding: the synthesis tool picks one, and
you influence it with an attribute. The \SV{} enumeration has an explicit
base type, \code{logic [1:0]}, and that is what goes into the flip-flops. The
\SV{} version is more honest about the hardware and less convenient when you
want the tool to choose one-hot for you.

The \VHDL{} \code{case} must cover every value of \code{state\_t} or the
analyser complains, and there are exactly three. The \SV{} \code{case} covers
three of four possible two-bit values, so it needs a \code{default}, and
\code{unique} asks the simulator to check at run time that exactly one branch
matched. \VHDL{} gets that check for free from its type system.

\code{always\_comb} is not quite \code{process (all)}. Both derive the
sensitivity list. \code{always\_comb} additionally executes once at time
zero, and it includes the contents of any function it calls in the
sensitivity list. Neither of them checks for an inferred latch: the default
assignments at the top of the block are your responsibility in both
languages.

\subsection{Registers and checks}

\begin{rosetta}
\begin{rleft}
p_regs : process (clk_i, rst_ni) is
begin
  if rst_ni = '0' then
    state_q <= ST_IDLE;
    x_q     <= (others => '0');
  elsif rising_edge(clk_i) then
    state_q <= state_d;
    x_q     <= x_d;
  end if;
end process p_regs;
\end{rleft}
\begin{rright}
always_ff @(posedge clk_i or
            negedge rst_ni) begin
  if (!rst_ni) begin
    state_q <= ST_IDLE;
    x_q     <= '0;
  end else begin
    state_q <= state_d;
    x_q     <= x_d;
  end
end
\end{rright}
\end{rosetta}

This is the pair that will convince you the two languages are the same
language wearing different clothes. \code{'0} is the one novelty: an unsized
fill literal that takes the width of its context, so it works for a two-bit
enumeration and a twenty-bit datapath without being told the width.

The checks are where \SV{} pulls ahead. Both designs carry an immediate
assertion on the input range:

\begin{rosetta}
\begin{rleft}
assert req_theta_i
         <= to_signed(PI_EXT, AW)
   and req_theta_i
         >= to_signed(-PI_EXT, AW)
  report "|theta| > pi"
  severity error;
\end{rleft}
\begin{rright}
assert (req_theta_i <=  angle_t'(PI_EXT) &&
        req_theta_i >= -angle_t'(PI_EXT))
  else $error("cordic_rot: |theta| > pi (%0d)",
              req_theta_i);
\end{rright}
\end{rosetta}

But the second check --- that a stalled request holds its payload steady --- is
a \emph{temporal} property, and there the two diverge completely. In \SV{} it
is four lines:

\begin{svcode}
property p_req_stable;
  @(posedge clk_i) disable iff (!rst_ni)
    (req_valid_i && !req_ready_o) |=>
      $stable(req_x_i) && $stable(req_y_i) && $stable(req_theta_i);
endproperty
a_req_stable : assert property (p_req_stable)
  else $error("cordic_rot: request payload changed while stalled");
\end{svcode}

In \VHDL{} it is a process with five variables that remembers the previous
cycle by hand, twenty lines in \cref{lst:case:vhdlrot}, or it is PSL, which
is a second language with patchy tool support. This is the clearest single
win for \SV{} in the whole design.

\clearpage
\svfile{code/cordic/sv/cordic_pkg.sv}{\file{cordic\_pkg.sv} --- the
SystemVerilog package.}{lst:case:svpkg}

\clearpage
\svfile{code/cordic/sv/cordic_rot.sv}{\file{cordic\_rot.sv} --- the
SystemVerilog implementation.}{lst:case:svrot}

\clearpage
\section{The Two Implementations Compared}
\label{sec:case:compare}

\Cref{tab:case:lines} counts the lines. Both files are commented at the same
density, so the ratio is meaningful.

\begin{table}[htbp]
  \centering
  \small
  \caption{Source size. Comments and blank lines are included, because the
  two files were written to the same commenting standard.}
  \label{tab:case:lines}
  \begin{tabular}{@{}lrrl@{}}
    \toprule
    & VHDL & SystemVerilog & \\
    \midrule
    package        & 148 & 140 & \\
    implementation & 253 & 242 & \\
    testbench      & 290 & 689 & (interface + package + top)\\
    \midrule
    design total   & 401 & 382 & \\
    \bottomrule
  \end{tabular}
\end{table}

The design is a wash: 401 lines against 382, five per cent. That is worth
saying plainly, because the folklore holds that \SV{} is dramatically
terser. For \emph{structural} RTL it is not. The savings that folklore is
remembering come from interfaces, from packed structs, and above all from
testbenches, and this design uses none of the first two.

The testbench is the opposite: 290 lines against 689, and the \SV{} version
does more. That is the real comparison, and \cref{sec:case:vhdltb} and
\cref{sec:case:svtb} put the two side by side.

Four places where the languages genuinely differ for this design:

\textbf{The type system did work for the \VHDL{}.} Three width errors were
caught by the analyser while this design was being written, before anything
was simulated. The equivalent \SV{} errors surfaced as \code{WIDTHEXPAND} and
\code{WIDTHTRUNC} warnings from \code{verilator --lint-only -Wall}, which is
a perfectly good substitute \emph{provided you run it}, and which no student
runs by default. The mitigation is a \code{make lint} target that fails the
build.

\textbf{The constant table is a draw with a portability asterisk.} Both
languages compute it at elaboration from a transcendental function. \VHDL{}'s
version works in every tool that implements the standard;
\SV{}'s does not (see the trap in \cref{sec:case:sv}).

\textbf{The enumeration encoding favours \VHDL{} for synthesis and \SV{} for
debugging.} A \VHDL{} state type lets the tool pick the encoding; a \SV{} one
tells you what is in the flip-flops. Which you want depends on whether you
are closing timing or reading a waveform.

\textbf{Assertions favour \SV{} decisively.} One temporal property is four
readable lines in \SV{} and twenty unreadable ones in \VHDL{}. Multiply that
by the ten protocol properties a real bus interface has.

\begin{notebox}{And one thing this design should have used}
\code{ieee.fixed\_pkg} was available for the \VHDL{} and was not used, so
that the two implementations would be structurally comparable. A production
\VHDL{} design of this kind should use it: it would have made
\cref{tab:case:formats} machine-checkable instead of documentation.
\end{notebox}

\clearpage
\section{The VHDL Testbench}
\label{sec:case:vhdltb}

\Cref{lst:case:vhdltb} is a self-checking \VHDL{} testbench, written the way a
\VHDL{} designer writes one: a stimulus process, a checking process,
procedures instead of classes, and a golden model built from
\code{ieee.math\_real}. It is the baseline the rest of the chapter is measured
against.

\subsection{The parts, and what each one costs}

\textbf{The generic is the \VHDL{} plusarg.}

\begin{vhdlcode}
entity tb_cordic_rot is
  generic (
    G_N_RANDOM : natural  := 300;
    G_TOL_LSB  : natural  := 12;
    G_SEED1    : positive := 7;
    G_SEED2    : positive := 13
  );
end entity tb_cordic_rot;
\end{vhdlcode}

Elaborate it with \code{vsim -gG\_N\_RANDOM=1000 tb\_cordic\_rot} and you
change the run without touching the source. Unlike a \SV{} plusarg, a generic
is fixed at elaboration rather than read at run time, so changing it means
elaborating again, not merely restarting.

\textbf{The clock, with a stop condition.}

\begin{vhdlcode}
clk <= not clk after CLK_PERIOD / 2 when not halt else '0';
\end{vhdlcode}

A \VHDL{} simulation traditionally ends by running out of events, so the
clock has to be stoppable. \code{std.env.finish} (VHDL-2008) is the modern
alternative and the testbench uses both: it sets \code{halt} and then calls
\code{finish}.

\textbf{Driving on the falling edge.} The \code{send} procedure asserts
\code{req\_valid} and then waits for falling edges:

\begin{vhdlcode}
procedure send (x, y, theta : integer) is
begin
  req_x     <= to_signed(x, DW);
  req_theta <= to_signed(theta, AW);
  req_valid <= '1';
  loop
    wait until falling_edge(clk);
    exit when req_ready = '1';
  end loop;
  wait until rising_edge(clk);
  wait until falling_edge(clk);
  req_valid <= '0';
end procedure send;
\end{vhdlcode}

This is the convention that keeps the testbench out of a race with the DUT,
and it works. It is also a lie about the hardware: no real driver changes its
outputs half a cycle before the clock. A \SV{} clocking block expresses the
truth --- drive $1$~ns after the edge, sample the value that existed just
before it --- and does not need the convention.

\textbf{No constraint solver.}

\begin{vhdlcode}
impure function rand_int (lo, hi : integer) return integer is
begin
  uniform(seed1, seed2, r);
  return lo + integer(floor(r * real(hi - lo + 1)));
end function rand_int;
\end{vhdlcode}

\code{uniform} from \code{ieee.math\_real} is the whole of \VHDL{}'s random
number support. A uniform distribution over an interval is six lines; a
distribution weighted toward the quadrant boundaries is however many lines you
care to write. In \SV{} the same thing is a \code{dist} clause inside a
\code{constraint} block. This is a real productivity difference, and it is the
one OSVVM exists to close.

\textbf{The structural point: monitor and scoreboard cannot be separated.}

\begin{vhdlcode}
p_check : process is
  variable pend   : req_array_t(0 to Q_DEPTH - 1);
  variable p_head : natural := 0;
  variable p_tail : natural := 0;
begin
  ...
  if req_valid = '1' and req_ready = '1' then       -- monitor
    pend(p_tail) := (x => to_integer(req_x), ...);
    p_tail := (p_tail + 1) mod Q_DEPTH;
  end if;
  if rsp_valid = '1' and rsp_ready = '1' then       -- scoreboard
    cur := pend(p_head);
    ...
\end{vhdlcode}

The observation of the request channel and the checking of the response
channel are in one process because they have to share the pending queue, and
a \VHDL{} process variable is the only cheap way to share mutable state. There
is no mailbox to put a transaction into and no dynamic object to put in it.

The consequence is not that the testbench is worse --- it works, and it found
bugs --- but that it does not decompose. You cannot lift the monitor out and
reuse it against a gate-level netlist, or run two of them, or subscribe a
coverage collector to the same stream, without rewriting.

\begin{gotcha}{Sharing a queue through a signal does not work}
The first version of this testbench logged each accepted request into a
\code{signal} array indexed by a counter, for the checking process to read.
Every entry read back as zero. A signal assignment inside a process takes
effect after the process suspends, the index is itself a signal, and the two
delta-cycle behaviours interact in a way that is correct by the standard and
useless in practice. The fix was to make the checking process observe the
request channel itself, which is what a \SV{} monitor does anyway.
\end{gotcha}

\subsection{Running it}

Three commands: make a library, analyse the three files in dependency order,
elaborate and run the top.

\begin{shcode}
vlib work
vcom -2008 cordic_pkg.vhd cordic_rot.vhd tb_cordic_rot.vhd
vsim -c -gG_N_RANDOM=300 tb_cordic_rot -do "run -all; quit"
\end{shcode}

Everything the testbench has to say about itself comes from \code{report}
statements, so the summary reads the same whatever tool executed it. Stripped
of the per-simulator prefix that decorates each line, the run prints:

\begin{txtcode}
=== cordic_rot testbench (VHDL) ===
ITER=14  random=300
------------------------------------------------
scoreboard: 308 transactions checked, 0 failed
worst error: 2.5458225901211335 LSB (tolerance 12)
------------------------------------------------
*** TEST PASSED ***
\end{txtcode}

308 transactions rather than 300, because the directed list runs first and the
random phase is counted on top of it. The worst error, $2.55$ LSB, is the
number \cref{sec:case:format} predicted from the error budget.

Before any of that, at time zero, the run also emits a handful of warnings
from \code{numeric\_std} itself, of the form
\code{NUMERIC\_STD.">": metavalue detected, returning FALSE}. They are worth
understanding rather than silencing by reflex.

\begin{notebox}{Those metavalue warnings are \VHDL{} telling you something}
At time zero, before the first assignment has propagated, the DUT's inputs
are \code{'U'}. The comparison \code{z\_ext > HALF\_PI\_S} on a vector
containing \code{'U'} cannot produce a meaningful answer, so
\code{numeric\_std} says so and returns \code{FALSE}. The message comes from
the IEEE library, not from the simulator, so every compliant tool reports it.
Verilog has no \code{'U'}: an uninitialised variable is \code{X}, and
\code{X > 5} quietly evaluates to \code{X} and then to false in an \code{if},
with no warning at all. \VHDL{} is being more informative, not more broken.

For a regression, where the noise buries the result, put

\begin{tclcode}
set NumericStdNoWarnings 1
\end{tclcode}

in the \code{do} file before \code{run -all}. Do not set it while you are
debugging a reset: those warnings are exactly the evidence that something is
still \code{'U'} when you thought it was driven.
\end{notebox}

\clearpage
\vhdlfile{code/cordic/vhdl/tb_cordic_rot.vhd}{\file{tb\_cordic\_rot.vhd} ---
the VHDL testbench, complete.}{lst:case:vhdltb}

\clearpage
\section{The SystemVerilog Testbench}
\label{sec:case:svtb}

Three files, 689 lines, and every structural idea from \refverificationbook in
working form: an interface with clocking blocks, six classes, mailboxes
between them, a directed phase followed by a random one, and an end of test
that drains rather than guesses.

\subsection{The interface}

\Cref{lst:case:svif} is the pin bundle. The part that matters is the clocking
block:

\begin{svcode}
clocking drv_cb @(posedge clk);
  default input #1step output #1ns;
  output req_valid, req_x, req_y, req_theta;
  output rsp_ready;
  input  req_ready, rsp_valid, rsp_x, rsp_y;
endclocking
\end{svcode}

\code{input \#1step} samples in the Preponed region: the value that existed
immediately before the clock edge, before any non-blocking update. \code{output
\#1ns} drives one nanosecond after the edge. Together they remove the entire
class of bug that the \VHDL{} testbench avoids by driving on the falling edge,
and they model what a real driver does.

\subsection{The six classes}

\Cref{lst:case:svtbpkg} is the testbench package. The shape is exactly
\refverificationbook:

\begin{itemize}
  \item \code{cordic\_txn} --- the transaction, with \code{rand} fields,
        constraints (including a \code{dist} that biases toward the quadrant
        boundaries), and \code{convert2string};
  \item \code{cordic\_drv} --- the driver, which knows the handshake and
        nothing else, and runs two threads: one feeding requests, one
        applying random back-pressure to the response channel;
  \item \code{cordic\_mon} --- the monitor, strictly passive, holding a queue
        of accepted-but-unanswered requests;
  \item \code{cordic\_sb} --- the scoreboard and its specification-level
        reference model;
  \item \code{cordic\_cov} --- the coverage collector, with a covergroup that
        crosses the angle quadrant against the sign of $x_0$;
  \item \code{cordic\_env} --- construction, connection, start-up.
\end{itemize}

The reference model is worth quoting because of what it does \emph{not} say:

\begin{svcode}
function automatic void predict(input cordic_txn t,
                                output real x_ref, output real y_ref);
  real th = q_to_real(t.theta, QA);
  real xr = q_to_real(t.x0,    QF);
  real yr = q_to_real(t.y0,    QF);
  x_ref = K_GAIN * (xr * $cos(th) - yr * $sin(th));
  y_ref = K_GAIN * (xr * $sin(th) + yr * $cos(th));
endfunction
\end{svcode}

No micro-rotations, no guard bits, no quadrant folding. It states what the
block is supposed to compute. Every implementation detail is something it can
therefore catch.

\subsection{The portability shim}

At the top of \file{tb\_cordic\_pkg.sv} sit four macros, and they need
explaining because they are the one piece of the design that exists for a
tool rather than for the hardware.

\begin{svcode}
`ifdef VERILATOR
  `define CB_DRV       vif
  `define CB_MON       vif
  `define EDGE_DRV     @(posedge vif.clk)
  `define EDGE_MON     @(posedge vif.clk)
`else
  `define CB_DRV       vif.drv_cb
  `define CB_MON       vif.mon_cb
  `define EDGE_DRV     @(vif.drv_cb)
  `define EDGE_MON     @(vif.mon_cb)
`endif
\end{svcode}

\tool{Verilator} 5.020 has two problems with clocking blocks reached through
a virtual interface, and the second one is nasty.

The first is honest: a \code{modport} that exports a clocking block is
rejected with \code{Unsupported: Modport clocking}. You see the message, you
work around it.

The second is not. A \emph{non-blocking} assignment through a virtual
interface to a signal that is a clocking-block output is silently dropped.
The write never happens; nothing is reported. This was not assumed --- it was
reduced to a ten-line test case and confirmed:

\begin{svcode}
interface bif (input logic clk);
  logic a, c;
`ifndef NOCB
  clocking cb @(posedge clk); output a; endclocking
`endif
endinterface
// ... a class with `virtual bif v` doing:
//        v.a <= 1'b1;   // in the clocking block
//        v.c <= 1'b1;   // not in any clocking block
\end{svcode}

\begin{txtcode}
=== WITH clocking block ===          === WITHOUT clocking block ===
[ 5] a=0 c=0                         [ 5] a=0 c=0
[15] a=0 c=1                         [15] a=1 c=1
[25] a=0 c=1                         [25] a=1 c=1
\end{txtcode}

\code{c} takes the assignment; \code{a} does not, purely because \code{a} is
named in a clocking block. The symptom in the real testbench was that
\code{req\_valid} never went high, the DUT never left \code{IDLE}, and the
run timed out after 42\,600 cycles with zero transactions completed --- with
no error message pointing anywhere near the cause.

\begin{gotcha}{A silently dropped write is the worst failure mode there is}
An unsupported construct that produces an error message costs you ten
minutes. One that produces wrong behaviour with no message costs you an
afternoon, and it teaches you to distrust the tool rather than the code. This
is the strongest argument in the book for having a second simulator: the same
sources under \tool{QuestaSim} and under \tool{Verilator} would have located
this in one run.
\end{gotcha}

The shim keeps one set of sources runnable in both worlds. Under any IEEE
1800 simulator the macros expand to the clocking block, with its sampling and
drive skew; under \tool{Verilator} they expand to the raw signals plus an
explicit clock edge, which loses the skew and keeps the test.

\subsection{Ending the test}

\begin{svcode}
task automatic wait_drain(int unsigned expected, int unsigned timeout_cyc);
  int unsigned waited = 0;
  while (sb.n_checked < expected && waited < timeout_cyc) begin
    `EDGE_MON;
    waited++;
  end
  if (sb.n_checked < expected)
    $fatal(1, "timeout: only %0d of %0d transactions completed",
           sb.n_checked, expected);
endtask
\end{svcode}

Count what went in, wait until that many have been checked, fail if they
never arrive. This is precisely what UVM's objection mechanism automates, and
writing it once by hand is the fastest way to understand why the mechanism
exists (\refverificationbook).

\subsection{Running it}

\begin{txtcode}
$ make sv
=== cordic_rot testbench ===============================
 ITER=14  DW=16.Q14  AW=16.Q13  random=300
--------------------------------------------------------
 scoreboard: 308 transactions checked, 0 failed
 worst error: 2.82 LSB (tolerance 12 LSB)
--------------------------------------------------------
 *** TEST PASSED ***
 functional coverage: not collected (Verilator build)
\end{txtcode}

Note the last line. The covergroups are compiled out under \tool{Verilator},
and so are the constraints: \code{cordic\_txn::do\_randomize()} falls back to
\code{\$urandom\_range}. Under \tool{QuestaSim} the same three files run with
constrained random stimulus and a coverage number at the end. That is what
\refsimulationbook means when it says \tool{Verilator} implements a subset:
the subset is large enough to run this testbench and small enough that two of
its most valuable features go quiet.

\clearpage
\svfile{code/cordic/sv/cordic_if.sv}{\file{cordic\_if.sv} --- the interface,
its clocking blocks and its modports.}{lst:case:svif}

\clearpage
\svfile{code/cordic/sv/tb_cordic_pkg.sv}{\file{tb\_cordic\_pkg.sv} --- the
layered class testbench, complete.}{lst:case:svtbpkg}

\clearpage
\svfile{code/cordic/sv/tb_cordic_top.sv}{\file{tb\_cordic\_top.sv} --- the
top level: clock, interface, DUT, test.}{lst:case:svtbtop}

\clearpage
\section{The C++ Testbenches}
\label{sec:case:cpp}

Three, in increasing order of structure. All three drive the same
\tool{Verilator}-compiled model of the \SV{} RTL.

\subsection{Direct: \file{tb\_cordic.cpp}}

\Cref{lst:case:cpptb} is the shortest thing that does a real job: a clock
loop, a handshake, a \code{<cmath>} reference and an error counter. Two
hundred transactions, a tolerance check, an exit code.

The part to read carefully is the cycle:

\begin{cppcode}
void tick() {
  dut_->clk_i = 0;
  dut_->eval();          // settle: pre-edge values are now valid
  ctx_->timeInc(5000);   // 5 ns at 1 ps precision
  dut_->clk_i = 1;
  dut_->eval();          // the rising edge
  ctx_->timeInc(5000);
  ++cycles_;
}
\end{cppcode}

Set the inputs while the clock is low, call \code{eval()} so the
combinational cone settles, and only then raise the clock. That sequence is
the C++ spelling of ``sample in the Preponed region''. Read \code{req\_ready\_o}
after the low-phase \code{eval()} and you get the value the DUT will act on at
the edge; read it after the high-phase \code{eval()} and you get next
cycle's, which is a bug you will spend an hour on.

\begin{gotcha}{\code{timeInc} counts in the context's time precision}
A \code{VerilatedContext} defaults to 1 ps precision, so a 5 ns half period is
\code{timeInc(5000)}, not \code{timeInc(5)}. Get it wrong and the logic is
unaffected --- nothing in the design depends on absolute time --- but every
timestamp in the log and in the VCD is out by a factor of a thousand, and the
waveform lines up with nothing.

This is not a hypothetical. In the sources shipped with this book,
\file{tb\_cordic.cpp} uses \code{timeInc(5)} while \file{tb\_vectors.cpp} and
\file{main\_uvmlike.cpp} use \code{timeInc(5000)}. The second pair is right.
The first is a real bug, found while writing this chapter and left in place
so that you can see what it looks like: nothing fails, nothing warns, and the
reported simulation time is wrong.
\end{gotcha}

\begin{txtcode}
$ make cpp
----------------------------------------------
 200 transactions, 0 failures, worst error 2.68 LSB
 3206 clock cycles simulated
 *** TEST PASSED ***
\end{txtcode}

\subsection{Bit-exact: \file{tb\_vectors.cpp}}

\Cref{lst:case:cppvec} does something the other testbenches cannot: it
compares against the Python model \emph{exactly}. The vectors arrive as a
generated header, so a stale table is a compile error rather than a mystery.

\begin{cppcode}
#include "cordic_vectors.h"        // generated by gen_vectors.py
...
if (got_x != v.x || got_y != v.y) { ... ++failures; }
\end{cppcode}

The other thing this file exists to demonstrate is sign extension:

\begin{cppcode}
static inline int16_t as_signed(uint16_t raw) {
  return static_cast<int16_t>(raw);
}
\end{cppcode}

\begin{gotcha}{Verilator ports carry no signedness}
A 16-bit port becomes a \code{uint16\_t} (\code{SData}). \code{logic signed [15:0]}
and \code{logic [15:0]} produce the same C++ type. If you read
\code{rsp\_x\_o} and compare it to a negative expected value without the cast,
every positive result matches and every negative one fails --- or, worse, you
compare two unsigned values and everything passes. This is the single most
common bug in a first C++ testbench.
\end{gotcha}

\begin{txtcode}
$ make bitexact
----------------------------------------------
 bit-exact check against the Python model
 vectors generated with seed 1
 315 checked, 0 mismatches, 5046 clock cycles
 *** TEST PASSED ***
----------------------------------------------
\end{txtcode}

A tolerance check and a bit-exact check fail on different things. The
tolerance check fails when the mathematics is wrong: a missing iteration, a
wrong sign, a bad table entry. The bit-exact check fails when a
\emph{rounding decision} has drifted: a shift that rounds toward zero instead
of toward $-\infty$, a guard bit dropped one stage early, a saturation limit
off by one. Neither subsumes the other, and a design of this kind wants both.

\subsection{UVM-shaped: \file{svtlm.h}}

\Refsimulationbook builds the framework; here is the CORDIC sitting in it.
The component tree is real, printed by the running program:

\begin{txtcode}
$ make uvmlike
[test.env.agt.seq      ] will generate 200 transactions
--- component tree ---
test (CordicTest)
  env (CordicEnv)
    agt (CordicAgent)
      seq (CordicSeq)
      drv (CordicDrv)
      mon (CordicMon)
    sb (CordicSb)
    cov (CordicCov)
--- run ---
--- report ---
[test.env.agt.drv      ] drove 200 transactions
[test.env.sb           ] 200 checked, 0 failed, worst 2.68 LSB (tol 12)
[test.env.cov          ] theta bins hit 4/5  [62 47 0 46 45]
*** TEST PASSED ***
\end{txtcode}

Compare that tree with the UVM hierarchy in \refverificationbook. It is the same
tree, with the same names, built by a factory from a string and configured
from a key-value store, in about two hundred lines of header-only C++ and no
library at all.

Look at the coverage line, which is the most interesting output in this
chapter: \code{[62 47 0 46 45]}. Five bins over the angle range, and the
middle one --- $\theta$ exactly zero --- was never hit in two hundred random
transactions. Of course it was not: the probability of a uniform draw over
51\,473 values landing on zero is $2\times10^{-5}$. That is what a coverage
hole looks like, and the fix is the directed list at the top of
\cref{lst:case:svtbtop}, which sends $\theta = 0$ deliberately.

\begin{ruleofthumb}
Random stimulus will not find the corner you can name. Name it, send it, and
use coverage to tell you about the ones you could not name.
\end{ruleofthumb}

\clearpage
\cppfile{code/cordic/verilator/tb_cordic.cpp}{\file{tb\_cordic.cpp} --- the
direct C++ testbench.}{lst:case:cpptb}

\clearpage
\cppfile{code/cordic/verilator/tb_vectors.cpp}{\file{tb\_vectors.cpp} --- the
bit-exact C++ testbench.}{lst:case:cppvec}

\clearpage
\cppfile{code/cordic/verilator/svtlm.h}{\file{svtlm.h} --- the UVM-shaped
framework: Component, AnalysisPort, Factory, ConfigDb.}{lst:case:svtlm}

\clearpage
\cppfile{code/cordic/verilator/cordic_env.h}{\file{cordic\_env.h} --- the
CORDIC components built on \file{svtlm.h}.}{lst:case:cppenv}

\clearpage
\cppfile{code/cordic/verilator/main_uvmlike.cpp}{\file{main\_uvmlike.cpp} ---
the test layer and the clock loop.}{lst:case:cppmain}

\clearpage
\section{The Python Testbench}
\label{sec:case:python}

\subsection{Generated everything}

\Cref{lst:case:gen} is one script with four outputs, and it is the practical
face of \refsimulationbook:

\begin{itemize}
  \item \file{vectors/cordic\_vectors.txt} --- a plain text vector file that
        any testbench can read;
  \item \file{verilator/cordic\_vectors.h} --- a C++ header with
        \code{constexpr} arrays, consumed by \file{tb\_vectors.cpp};
  \item \file{sv/cordic\_vectors.svh} --- the same table as \SV{}
        \code{localparam} arrays;
  \item \file{sv/cordic\_atan\_rom.svh} --- the arctangent ROM as generated
        source, for the tools that cannot compute it themselves.
\end{itemize}

The vector file carries its own provenance, which matters more than it looks:

\begin{txtcode}
# cordic_rot test vectors
# generated 2026-09-05 06:19:35
# generator  gen_vectors.py  seed=1
# format     DW=16 QF=14 AW=16 QA=13 N=14 GL=2 GM=2
# count      315
# x0    y0    theta xexp  yexp
26dd  0000  0000  3fff  0001
26dd  0000  6488  c000  0001
\end{txtcode}

Without the header, a vector file is an unreproducible blob that happens to
be readable. With it, anyone can regenerate it and see whether it changed.

\begin{gotcha}{Comment characters depend on the reader}
This file uses \code{\#} because it is read by \code{\$fgets}/\code{\$sscanf},
by \code{std.textio} and by a C++ \code{ifstream}, all of which are happy to
skip a line you tell them to skip. A file destined for \code{\$readmemh} is
different: \code{\$readmemh} accepts only \code{//} and \code{/* */}
comments, and a \code{\#} in it is a syntax error --- or, in some tools,
silently the value zero. Decide who reads the file before you choose the
comment character.
\end{gotcha}

\subsection{cocotb: the same test against both languages}

\Cref{lst:case:cocotb} is a testbench with no HDL in it at all. The top level
of the simulation is the DUT; \file{cordic\_model.py} is imported directly and
used as a bit-exact reference; the driver, the monitor and the scoreboard are
coroutines.

\begin{pycode}
async def monitor(dut, queue, results, expected):
    while len(results) < expected:
        await RisingEdge(dut.clk_i)
        if dut.rsp_valid_o.value == 1 and dut.rsp_ready_i.value == 1:
            txn = queue.popleft()
            txn.x = dut.rsp_x_o.value.to_signed()
            txn.y = dut.rsp_y_o.value.to_signed()
            results.append(txn)
\end{pycode}

\begin{gotcha}{\code{.value} is unsigned}
\code{int(dut.rsp\_x\_o.value)} on the result $-16384$ gives $49152$. The
test \code{test\_value\_semantics} in \cref{lst:case:cocotb} exists purely to
print all four spellings side by side so that you never have to rediscover
this. In cocotb 2.x the signed conversion is \code{.to\_signed()}; in 1.x it
was \code{.signed\_integer}. The equivalent 1.x spellings for the rest are
\code{units=} rather than \code{unit=} on \code{Clock}, and \code{.binstr}
rather than \code{str(v)}.
\end{gotcha}

The headline is the Makefile. Two variables choose the language and the
backend; nothing else in the flow moves:

\begin{makecode}
ifeq ($(SV),1)
  SIM             := verilator
  TOPLEVEL_LANG   := verilog
  VERILOG_SOURCES := $(ROOT)/sv/cordic_pkg.sv $(ROOT)/sv/cordic_rot.sv
else
  SIM             := questa
  TOPLEVEL_LANG   := vhdl
  VHDL_SOURCES    := $(ROOT)/vhdl/cordic_pkg.vhd $(ROOT)/vhdl/cordic_rot.vhd
  COMPILE_ARGS    += -2008
endif
\end{makecode}

One Python file, two languages, two backends, no change to the test.
\code{SIM = questa} reaches the VHDL through the FLI, the C interface
\tool{QuestaSim} exposes for exactly this purpose; \code{SIM = verilator}
reaches the \SV{} through VPI. \code{COMPILE\_ARGS} is passed to the analysis
step, which is where a language-revision switch belongs. Put it in
\code{SIM\_ARGS} instead and it goes to the run, which does not understand it.

cocotb prints its own summary table, and it prints it identically whichever
backend ran underneath:

\begin{txtcode}
$ make cocotb
*******************************************************************
** TEST                             STATUS  SIM TIME (ns)  RATIO  **
*******************************************************************
** test_cordic.test_directed          PASS        1820.00  121098 **
** test_cordic.test_random            PASS       16070.00  128429 **
** test_cordic.test_sweep             PASS       29030.00  128453 **
** test_cordic.test_value_semantics   PASS         230.00   98399 **
*******************************************************************
** TESTS=4 PASS=4 FAIL=0 SKIP=0                  47150.00  127368 **
*******************************************************************
\end{txtcode}

Four tests, four passes, and the third of them is a 181-point sweep of the
full circle checked bit for bit against \file{cordic\_model.py}.

Two honest limitations, both confirmed on the machine this book was built on:

\begin{itemize}
  \item \textbf{The \SV{} run needs a newer \tool{Verilator}.} cocotb 2.1.0
        refuses anything before 5.036: \code{cocotb requires Verilator 5.036
        or later, but using 5.020. Stop.} The check is on the \tool{Verilator}
        backend only, and it is a version check rather than a missing feature:
        upgrade and the same test runs.
  \item \textbf{\tool{Icarus Verilog} cannot compile the package at all},
        because of the constant-function ROM. That is not a cocotb problem,
        and \file{cordic\_atan\_rom.svh} is the way around it.
\end{itemize}

\clearpage
\pyfile{code/cordic/python/gen_vectors.py}{\file{gen\_vectors.py} --- vectors,
a C++ header, an \SV{} include and the ROM, from one model.}{lst:case:gen}

\clearpage
\pyfile{code/cordic/python/test_cordic.py}{\file{test\_cordic.py} --- the
cocotb testbench. The same file drives the VHDL and the
SystemVerilog.}{lst:case:cocotb}

\clearpage
\section{Build and Run}
\label{sec:case:build}

\Cref{lst:case:make} is the whole build. Six ways of exercising one design,
and \code{make all} runs the five that need only free tools. The two \VHDL{}
targets, \code{make vhdl} and \code{make cocotb}, are separate because they
need a \VHDL{} simulator; the command lines in the Makefile are
\tool{QuestaSim}'s.

\begin{table}[htbp]
  \centering
  \small
  \caption{What checks what. ``The VHDL flow'' is \code{vcom}/\code{vsim} on
  a commercial simulator, as \cref{sec:case:vhdltb} sets it out. The last two
  rows are what tie the two implementations together.}
  \label{tab:case:checks}
  \begin{tabular}{@{}lllll@{}}
    \toprule
    Testbench & Lang. & Simulator & Reference & Check\\
    \midrule
    \file{tb\_cordic\_rot.vhd} & VHDL   & the VHDL flow & \code{ieee.math\_real} & tolerance\\
    \file{tb\_cordic\_top.sv}  & SV     & Verilator & \code{\$sin}/\code{\$cos} & tolerance\\
    \file{tb\_cordic.cpp}      & C++    & Verilator & \code{<cmath>}          & tolerance\\
    \file{main\_uvmlike.cpp}   & C++    & Verilator & \code{<cmath>}          & tolerance\\
    \file{tb\_vectors.cpp}     & C++    & Verilator & \code{cordic\_model.py} & bit-exact\\
    \file{test\_cordic.py}     & Python & the VHDL flow & \code{cordic\_model.py} & bit-exact\\
    \bottomrule
  \end{tabular}
\end{table}

The verification argument is worth stating explicitly, because it is the
reason for the last two rows:

\begin{enumerate}
  \item The \SV{} implementation is bit-exact against the Python model:
        \code{make bitexact}, 315 vectors, 0 mismatches.
  \item The \VHDL{} implementation is bit-exact against the \emph{same} Python
        model: \code{make cocotb}, four tests, 0 mismatches, including a
        181-point sweep of the full circle.
  \item Therefore the two implementations agree with each other, bit for bit,
        on every vector either of them has seen.
\end{enumerate}

That is a stronger statement than ``both testbenches pass'', and it costs one
extra script. Note also what the argument does \emph{not} depend on. Neither
step appeals to a particular simulator: step~1 needs some tool that runs the
\SV{}, step~2 needs some tool that runs the \VHDL{}, and the transitive
conclusion follows from the two comparisons against the one model. Change
either backend and the chain still closes.

There is one more target, and it produces something to look at rather than a
pass or a fail. \code{make waves} rebuilds the \SV{} testbench with tracing
enabled and runs a short twenty-transaction stimulus, leaving
\file{tb\_cordic.vcd} behind. A VCD is a file, not a viewer: open it in
\tool{GTKWave} or in \tool{Surfer}, or read it into a commercial simulator's
own wave window. \Refsimulationbook covers tracing properly, including why
you should name the scope you dump rather than dumping the whole hierarchy.

\begin{tipbox}{Lint first, every time}
\code{make lint} is \code{verilator --lint-only -Wall} over the synthesisable
sources, and it fails the build. In \VHDL{} the analyser does this job for
free. In \SV{} it is a target you have to write and a habit you have to
acquire, and it is the single highest-value line in the Makefile.
\end{tipbox}

\clearpage
\makefile{code/cordic/Makefile}{The build. \code{make all} runs lint, the
SystemVerilog class testbench, both C++ testbenches and the bit-exact check;
\code{make vhdl} and \code{make cocotb} are separate because they need a
\VHDL{} simulator.}{lst:case:make}

\section{What Is Not Verified}
\label{sec:case:gaps}

Six testbenches and a green regression are not the same thing as a verified
design. What has been established is that the \RTL{} computes the right
arithmetic on the vectors it has seen. Here is what has not, roughly in the
order a real project would address it.

\begin{enumerate}
  \item \textbf{Coverage is not closed.} The covergroups do not run under
        \tool{Verilator}, so the only coverage number in this chapter is the
        hand-rolled one in the C++ environment, and it reports a hole. Nobody
        has measured which states, which transitions or which angle ranges the
        300 random transactions actually exercised. Run the \SV{} testbench
        under a simulator with covergroups and close the holes; that is the
        first thing to do.

  \item \textbf{The saturating output is barely tested.} Two directed vectors
        reach it. There is no test that walks the input magnitude across the
        saturation boundary, and none at all for the negative limit, where an
        off-by-one is most likely.

  \item \textbf{Reset is tested once.} Reset is asserted at time zero and
        released. There is no test that asserts it mid-transaction, none that
        checks recovery, and none that varies its release relative to the
        clock edge.

  \item \textbf{The protocol is under-checked.} One concurrent assertion
        covers request stability. Nothing checks that \code{rsp\_valid\_o}
        stays asserted until \code{rsp\_ready\_i}, that the response payload
        is stable while stalled, or that no response appears without a
        request. Those are three more \code{assert property} statements and
        they belong in a \code{bind}-attached checker.

  \item \textbf{No gate-level simulation.} The design has not been
        synthesised. \tool{Verilator} is two-state, so a reset bug that
        depends on \code{X} propagation would not show up here at all --- see
        the \code{--x-assign} discussion in \refsimulationbook.

  \item \textbf{No formal.} The handshake properties are exactly the kind of
        thing a bounded model checker proves in seconds. \code{assume} the
        input protocol, \code{assert} the output protocol, and you get a
        proof rather than 308 samples.

  \item \textbf{One parameterisation.} Everything was run at
        \code{ITER = 14}. The parameter exists; nothing sweeps it.
\end{enumerate}

\begin{ruleofthumb}
A regression that passes tells you what you tested. Write down what you did
not.
\end{ruleofthumb}

\section{Exercises}
\label{sec:case:exercises}

\begin{enumerate}
  \item \textbf{Sweep \code{ITER}.} Run the testbench at \code{ITER} = 8, 10,
        12, 14 and 16 and plot the worst error. Explain the floor.
        \emph{Teaches: where the error budget actually comes from.}

  \item \textbf{Break it on purpose.} Change \code{>>>} to \code{>>} in
        \file{cordic\_rot.sv} and run every testbench. Which ones fail, how
        loudly, and how long does each take to point you at the line?
        \emph{Teaches: what a scoreboard is for.}

  \item \textbf{Make it build under Icarus.} Add a \code{`ifdef} to
        \file{cordic\_pkg.sv} that uses the generated
        \file{cordic\_atan\_rom.svh} instead of the constant function, and
        get \code{iverilog -g2012} to compile the design.
        \emph{Teaches: generated source, and why portability costs something.}

  \item \textbf{Add the missing protocol checks.} Write the three assertions
        listed in \cref{sec:case:gaps} in a separate checker module and attach
        it with \code{bind}, without editing \file{cordic\_rot.sv}.
        \emph{Teaches: \code{bind}, and assertion-based verification of code
        you do not own.}

  \item \textbf{Close the coverage.} Add a covergroup that crosses the angle
        quadrant with the sign of $y_0$ and with whether the output
        saturated, and add stimulus until it closes.
        \emph{Teaches: the coverage-driven loop.}

  \item \textbf{Vectoring mode.} Change the sign decision from
        $d_i = \operatorname{sign}(z_i)$ to
        $d_i = -\operatorname{sign}(y_i)$ and the block computes
        $\sqrt{x^2+y^2}$ and $\arctan(y/x)$ instead. Update the model first,
        then the \RTL{}.
        \emph{Teaches: that the golden model is the specification.}

  \item \textbf{Pipeline it.} Unroll the loop into $N$ stages so the block
        accepts a transaction every cycle. Work out what happens to
        \code{req\_ready\_o} and to \code{rsp\_valid\_o}, and what the monitor
        has to do differently.
        \emph{Teaches: that the interesting part of a pipeline is the
        handshake, not the datapath.}

  \item \textbf{Rewrite the testbench in UVM.} Take
        \file{tb\_cordic\_pkg.sv} and map each class onto its UVM base, using
        the table in \refverificationbook. Count the lines before and after.
        \emph{Teaches: exactly what UVM adds, and what it costs.}
\end{enumerate}

\section{Checklist}
\label{sec:case:checklist}

\begin{itemize}
  \item Write the number format down before the \RTL{}, and put it in the
        package. \SV{} will not check it for you.
  \item Write the golden model before the \RTL{}, from the specification, and
        never by reading the \RTL{}.
  \item Make the model bit-exact if the design is fixed-point. A tolerance
        check and a bit-exact check catch different bugs; run both.
  \item Compute constant tables at elaboration where the tools allow it, and
        generate them from a script where they do not. Put a
        \code{GENERATED} banner in the first line of anything generated.
  \item Put \code{make lint} first in the build and make it fail the build.
        It is the replacement for \VHDL{}'s analyser.
  \item Use \code{>>>} on signed values. Check every shift in the design
        before you check anything else.
  \item Reach the DUT through a clocking block. Where a tool cannot, say so
        in a comment at the top of the file and keep one set of sources.
  \item End the test by draining and timing out, never by waiting a fixed
        number of cycles.
  \item Send the corners you can name; use coverage for the ones you cannot.
  \item Run two simulators if you can. The bugs that cost a day are the ones
        one tool implements silently wrong.
  \item Write down what you did not verify, in the same commit as the
        testbench that passes.
\end{itemize}
